\documentclass[11pt]{article}
\usepackage[a4paper,margin=23mm]{geometry}
\usepackage[T1]{fontenc}
\usepackage{lmodern,amsmath,amssymb,amsthm,mathtools,microtype,booktabs}
\usepackage[colorlinks=true,linkcolor=blue,citecolor=blue,urlcolor=blue]{hyperref}
\newcommand{\E}{\mathbb E}
\newcommand{\Tr}{\operatorname{Tr}}
\newcommand{\tr}{\operatorname{tr}}
\newcommand{\HH}{\mathcal H}
\newcommand{\dd}{\mathrm d}
\newcommand{\norm}[1]{\left\lVert#1\right\rVert}
\newtheorem{theorem}{Theorem}
\title{Uniform Slater-orbit densities for linear selected preparations\\
\large Exact energy pinching in the one-excitation sector}
\author{Alan Day \and Chengheng\thanks{AI research assistant (OpenAI).}}
\date{8 November 2026\quad Local research draft}
\begin{document}\maketitle
\begin{abstract}
We prove a finite-dimensional, size-independent moment bound for the
orbit density of a linear Slater-adapted preparation. The key estimate
controls both variance operators of a rectangular Gaussian synthesis
map. Its input covariance is bounded using the Weyl orbit of the
Slater weight, with representation multiplicities retained. Standard
matrix-Gaussian integration by parts then bounds every fixed density
moment. For the normalized selected creation operator the second
moment is at most three, uniformly in the microscopic parameters.
This transfers exact energy-pinching errors from the covariant trace
to the one-selected excited state and closes its actual neutral-slope
correlators. Nonlinear selected preparations are left for a separate
estimate.
\end{abstract}

\section{Finite-$N$ setting and the result}
Let $\HH_m=\bigwedge^m\mathbb C^N$, $d_m=\binom Nm$ and $|O\rangle$
be a fixed Slater determinant. Use the original overlapping complex
Hermitian Gaussian $p$-body couplings. Express their GUE law in real
independent standard Gaussian coordinates $g_a$. Physical $U(N)$
rotations act on these coordinates by an orthogonal representation
$R(U)$ and on $\HH_m$ by $U_m=\bigwedge^m U$.

Take a homogeneous linear preparation $P_O(J)$ satisfying
\begin{equation}\label{eq:stab}
 P_O(k\cdot J)=k_m P_O(J)k_m^\dagger,
 \qquad k\in U(m)\times U(N-m),
 \quad a_N=\E\norm{P_O(J)|O\rangle}^2.
\end{equation}
The subgroup preserves the occupied plane. Normalized $M_N^\dagger$,
$h_N$, $u_N$ and fixed linear combinations obey this condition.
For $M_N^\dagger$, $a_N=1$ exactly.

\begin{theorem}\label{thm:main}
There is a positive covariant random matrix $\mathcal R_{P,N}$ such
that, for every covariant $Y_N(J)$ with $\E\tr_m(Y_N^\dagger Y_N)<\infty$,
\begin{equation}
 \E\langle O|P_O^\dagger Y_NP_O|O\rangle
    =\E\tr_m(\mathcal R_{P,N}Y_N),\qquad
 \E\tr_m\mathcal R_{P,N}^{,r}\le(2r-1)!!\,a_N^r
\end{equation}
for every integer $r\ge1$. Here $\tr_m=\Tr/d_m$.
\end{theorem}
This theorem is finite $N$; it does not assume double scaling or
replace finite-$N$ spectral projectors by their limiting functions.

\newpage
\section{The orbit synthesis map and its output covariance}
Haar measure on $G=U(N)$ is normalized to one. Put
\begin{equation}
 q_U(J)=U_m^\dagger P_O(U\cdot J)|O\rangle,\qquad
 Z(J)f=\sqrt{d_m}\int_G q_U(J)f(U)\,\dd U.
\end{equation}
The map $Z:L^2(G)\to\HH_m$ is a finite Gaussian matrix series after
restricting its domain to the finite-dimensional matrix-coefficient
space that supports all its entries. Define
\begin{equation}\label{eq:density}
 \mathcal R_{P,N}(J)=Z(J)Z(J)^\dagger
     =d_m\int_G |q_U(J)\rangle\langle q_U(J)|\,\dd U.
\end{equation}
Changing the Haar variable proves covariance of this matrix. Rotating
the Gaussian variable and using $Y(U\cdot J)=U_mY(J)U_m^\dagger$
proves the exact pairing identity in Theorem~\ref{thm:main}.

For the output covariance, Gaussian invariance first removes the
rotation from $P_O(U\cdot J)$ under expectation. Schur averaging
in the irreducible exterior-power representation then gives
\begin{equation}\label{eq:output}
 \E ZZ^\dagger
 =d_m\int_G U_m^\dagger
    \E[P_O(J)|O\rangle\langle O|P_O(J)^\dagger]U_m\,\dd U
 =a_N I_m.
\end{equation}
Controlling only \eqref{eq:output} would not suffice: a Gaussian
column $g\in\mathbb C^d$ has $\E gg^\dagger=I$ but
$\E\tr(gg^\dagger)^2=d+1$. The input variance must also be bounded.

Write $q_U=\sum_a g_a f_a(U)$ and introduce its coefficient vector
\begin{equation}
 \Psi_U=\sum_a e_a\otimes f_a(U)
     =[R(U)^T\otimes U_m^\dagger]\Psi
     =\rho(U^{-1})\Psi,\qquad
 \rho=R\otimes\bigwedge^m,\quad \norm\Psi^2=a_N.
\end{equation}
The real Hermitian noise basis makes $R(U)^T=R(U^{-1})$.
In ordered complex $J_{AB}$ labels the equivalent noise factor is
the conjugate adjoint action, so its torus weight is $B-A$;
the physical hop state has weight $O+A-B$.

For a diagonal unitary $D$, equation \eqref{eq:stab} gives
$q_D=\chi_O(D)^*q_I$, where $\chi_O(D)=\prod_{i\in O}D_{ii}$.
Consequently
\begin{equation}\label{eq:weight}
 \rho(D)\Psi=\chi_O(D)\Psi.
\end{equation}
The noise and physical weights therefore sum to the Slater weight.
This exact selection rule is what controls the other variance.

\newpage
\section{The Slater weight bounds the input covariance}
Decompose the finite-dimensional unitary representation as
$\bigoplus_\alpha V_\alpha\otimes\mathcal M_\alpha$ and let
$\Psi_\alpha$ be the component of $\Psi$. Equation \eqref{eq:weight}
holds in each nonzero component. The corresponding irreducible
representation $V_\alpha$ therefore contains the weight
$(1^m,0^{N-m})$.

Permutation matrices take that weight to every choice of $m$ occupied
coordinates. These $d_m$ weights are distinct, and their weight
spaces are mutually orthogonal. Hence
\begin{equation}
 d_\alpha:=\dim V_\alpha\ge d_m
 \qquad\text{whenever }\Psi_\alpha\ne0.
\end{equation}
No highest-weight assumption or multiplicity-free decomposition is
needed. Schur averaging, with multiplicity spaces retained, gives
\begin{align}
 F&:=d_m\int_G|\Psi_U\rangle\langle\Psi_U|\,\dd U\\
 &=\bigoplus_\alpha\frac{d_m}{d_\alpha}I_{V_\alpha}
       \otimes\sigma_\alpha,\qquad
 \sigma_\alpha=\Tr_{V_\alpha}
                   |\Psi_\alpha\rangle\langle\Psi_\alpha|.
\end{align}
Since $\sigma_\alpha\ge0$ and
$\norm{\sigma_\alpha}\le\Tr\sigma_\alpha
=\norm{\Psi_\alpha}^2\le a_N$, it follows that $0\le F\le a_N I$.

Let $Af=\sqrt{d_m}\int_G\Psi_U f(U)\dd U$.
Then $AA^\dagger=F$, while the kernel of $A^\dagger A$ is
$d_m\langle\Psi_U,\Psi_V\rangle$.
Gaussian covariance shows that this is precisely the kernel of
$\E Z^\dagger Z$. The nonzero spectra of $AA^\dagger$ and
$A^\dagger A$ coincide, so
\begin{equation}\label{eq:input}
 0\le\E Z^\dagger Z\le a_N I,
 \qquad \Tr(\E Z^\dagger Z)=d_m a_N.
\end{equation}
The bounds \eqref{eq:output} and \eqref{eq:input} are independent of
$N,p,m$ except through the preparation norm $a_N$.

For the selected creation block, this norm is one: it creates the
$L_N=\binom mp\binom{N-m}p$ orthogonal $p$-hop states from $O$,
with covariance $1/L_N$ per label. This also gives an elementary
check on \eqref{eq:output}: the unrotated output covariance is the
$p$-excitation projector divided by $L_N$, whose Haar average is
$I_m/d_m$.

\newpage
\section{Density moments and transfer of small covariant errors}
Hermitize the rectangular series:
\begin{equation}
 S=\begin{pmatrix}0&Z\\Z^\dagger&0\end{pmatrix}
   =\sum_a g_a A_a,\qquad
 V=\E S^2=\begin{pmatrix}a_N I_m&0\\0&\E Z^\dagger Z\end{pmatrix}.
\end{equation}
We have $V\le a_NI$ and $\Tr V=2d_m a_N$. For $r\ge2$, the
standard Gaussian integration-by-parts identity \cite[Sec.~6]{Tropp}
and the derivative of a matrix power give
\begin{equation}
 \E\Tr S^{2r}
 =\sum_a\sum_{k=0}^{2r-2}
     \E\Tr(A_aS^kA_aS^{2r-2-k}).
\end{equation}
In the eigenbasis of $S$, each summand is
$\sum_{ij}|(A_a)_{ij}|^2 x_j^k x_i^{2r-2-k}$.
Taking absolute values and applying weighted arithmetic--geometric
mean bounds it by $\Tr(A_a^2 S^{2r-2})$. This is the usual
matrix-Gaussian moment argument \cite[Sec.~7]{Tropp}. Thus
\begin{equation}
 \E\Tr S^{2r}\le(2r-1)\E\Tr(VS^{2r-2})
     \le(2r-1)a_N\E\Tr S^{2r-2}.
\end{equation}
Start at $\E\Tr S^2=\Tr V$, rather than at the ambient dimension.
Since $\Tr S^{2r}=2\Tr(ZZ^\dagger)^r$, induction proves
\begin{equation}\label{eq:moment}
 \E\tr_m\mathcal R_{P,N}^{,r}\le(2r-1)!!a_N^r.
\end{equation}
If $a_N=0$, the statement is immediate. The general Gaussian moment
method is established literature; the new input here is the orbit
variance estimate needed to apply it uniformly.

Let $\norm X_{r,m}=(\E\tr_m|X|^r)^{1/r}$. H\"older on the joint
disorder/matrix trace space now gives
\begin{align}\label{eq:holder}
 |\E\langle O|P_O^\dagger YP_O|O\rangle|
   &\le\sqrt3\,a_N\norm Y_{2,m},\\
 \E\norm{X P_O|O\rangle}^2
   &\le\sqrt3\,a_N\norm X_{4,m}^2
\end{align}
for covariant $Y,X$, respectively. Positivity of $Y$ is not required
in the first bound.

For a covariant post-preparation $B(J)$ there is also the exact
identity $\mathcal R_{BP}=B\mathcal R_PB^\dagger$. Consequently
\begin{equation}
 \norm{\mathcal R_{BP}}_{r,m}
 \le\norm B_{4r,m}^2\norm{\mathcal R_P}_{2r,m}.
\end{equation}
Fixed covariant polynomials have the required high moments; actual
unitary phases preserve them. The preparation order is $BP|O\rangle$.
It has not been interchanged with $PB|O\rangle$.

\newpage
\section{The actual neutral slope with exact energy pinching}
Return to half filling and the double-scaling limit
$\lambda_N=p^2/N\to\lambda>0$. In the notation of the preceding
stage, $h_N=H_N/\sqrt{C_m}$ and
\begin{equation}
 t_N=\frac{m\epsilon\alpha[\Phi_m(H_N)-\bar\kappa_m H_N]}
 {8\lambda_N C_m^{3/2}q_N\sqrt{\lambda_N}},\qquad
 b_N=(1-\mathcal P_{H_N})t_N.
\end{equation}
This $b_N$ is the normalized actual neutral slope. The pinching
$\mathcal P_{H_N}$ preserves whole degenerate energy blocks.
Time is dimensionless, conjugate to $h_N$; physical string energy
is $(J_s/\sqrt\lambda)h$.

The preceding central-sector argument already proved
\begin{equation}
 X_N:=\mathcal P_{H_N}t_N,
 \qquad \norm{X_N}_{r,m}\longrightarrow0
 \quad\text{for every fixed finite }r.
\end{equation}
It uses positivity of the two-energy measure, zero diagonal mass in
its $q$-Mehler limit, Portmanteau, high moments and pinching
contractivity. It does not interchange infinite time with $N$.
Applying \eqref{eq:holder} with $P=M_N^\dagger$, $a_N=1$, gives
\begin{equation}\label{eq:pinch}
 \E\norm{X_N M_N^\dagger|O\rangle}^2
      \le\sqrt3\norm{X_N}_{4,m}^2\longrightarrow0.
\end{equation}
This is the missing finite-$N$ estimate for the one-selected state.

For any fixed word $W$ of covariant insertions, telescope
$W(h_N,b_N)-W(h_N,t_N)$ and use tracial H\"older with sufficiently
high fixed moments. Its $S_2$ norm tends to zero. The first bound in
\eqref{eq:holder} transfers that difference to the selected state.
The already proved unpinched same-$J$ limit then yields
\begin{equation}\label{eq:limit}
 \lim_N\varphi_N\bigl[M_N W(h_N,b_N)M_N^\dagger\bigr]
   =\tau\bigl[\rho_1 W(h,u)\bigr],
\end{equation}
where $\tau$ is the standard two-color $q$-Gaussian trace and
\begin{equation}
 \rho_1=(\gamma h-\sqrt\lambda\gamma u)^2
          +[1-\sqrt q(1+\lambda)]I,\qquad \gamma=q^{1/4}.
\end{equation}
Any finite number of slope insertions and fixed actual $h_N$ time
phases inside the preparation sandwich are allowed. In particular,
\begin{equation}
 \lim_N\varphi_N(M_N b_N^2 M_N^\dagger)
      =1+\lambda\sqrt q(1+q).
\end{equation}
The finite-$N$ operation here is exact pinching, followed by $N\to
\infty$. It is stronger than the previous ordered finite-window
statement for this preparation class.

\newpage
\section{Exact finite checks, scope and the next interface}
Five finite Slater fixtures check 59 noise-plus-physical weight
identities and 122 permutation-frame counts. Two exact stabilizer
matrices for each of $(N,m,p)=(4,2,1),(4,2,2)$ check 624 entries,
including conjugation of the noise action and CAR signs.

An independent projective case provides the complete density. For
$m=p=1$, let $J$ be the single-particle Hermitian GUE matrix. Since
$q_\psi=(I-|\psi\rangle\langle\psi|)J|\psi\rangle/\sqrt{N-1}$,
Haar integration gives
\begin{equation}\label{eq:projective}
 \mathcal R(J)=\frac{NJ^2-2\Tr(J)J+
       [(\Tr J)^2+\Tr J^2]I/(N+1)}{(N-1)(N+2)}.
\end{equation}
The checker independently expands the three projector factors using
\begin{equation}
 \E_\psi\prod_{\ell=1}^k (|\psi\rangle\langle\psi|)_{i_\ell j_\ell}
 =\frac{\sum_{\pi\in S_k}\prod_\ell\delta_{i_\ell,j_{\pi(\ell)}}}
             {N(N+1)\cdots(N+k-1)},\qquad k\le3.
\end{equation}
This formula follows by twirling onto the symmetric tensor power.
Two exact complex fixtures at each $N=2,3,4$ verify all 58 entries
of \eqref{eq:projective} and its principal-minor positivity.
Only $N=2$ is half filled in this diagnostic family.

Independent Gaussian-coordinate Wick contraction at $N=2,3,4,5$
checks the mean $I$ and second moment with 978 coordinate pairs.
Writing $D=(N-1)(N+2)$, the result is
\begin{equation}
 \E\tr\mathcal R^2=
 \frac{2N^4-5N^2+3+\frac{2N+3}{(N+1)^2}(N^4-1)}{D^2}\le3.
\end{equation}
At $N=2$, $\mathcal R=(\chi_3^2/3)I$; its exact moments
$\prod_{j=0}^{r-1}(3+2j)/3^r$ obey \eqref{eq:moment}. Orders one
through eight are saved. No Haar Monte Carlo or asymptotic fit is used.

These checks supplement the general proof. The first checker attempt
compared differently simplified complex expressions; entrywise exact
simplification resolved the failed assertion. The initial diagnostic
is retained. Successful scientific output is sealed, with source,
versions, PDF build and renders bound by a separate SHA256 manifest.

The result covers linear preparations and covariant post-dressing,
including all covariant correlators in the $M_N^\dagger|O\rangle$
state. For multiple interlaced selected insertions, or a preparation
$M_N^\dagger B_N|O\rangle$, the synthesis map is generally nonlinear
Gaussian and the present moment recurrence does not apply.
The next task is that larger preparation class and the $J_2$/complete
bilocal interface, followed by the physical spatial, mass and
finite-$G$ dictionary.

\begin{thebibliography}{9}\small
\bibitem{Tropp} J. A. Tropp, \emph{Second-Order Matrix Concentration Inequalities},
\href{https://arxiv.org/abs/1504.05919}{arXiv:1504.05919};
\href{https://tropp.caltech.edu/papers/Tro16-Second-Order-Matrix-preprint.pdf}{author's preprint},
Secs.~6--7, Lemma 6.1 and Propositions 7.1, 7.3.
\end{thebibliography}
\end{document}

